% !TEX root = ../main.tex
\chapter{Introduction}
\label{sec:introduction}

\section{Motivation}
\label{sec:introduction-motivation}
Large language models are neural networks trained to predict the next word of a text, and the largest have trillions of parameters. Most of the leading open models of 2026, such as DeepSeek-V4, Kimi K3, GLM-5.3, Qwen3.8-Flash-Next, and MiMo-V2.6, are based on the \emph{mixture-of-experts} (MoE) architecture \cite{deepseek2026v4, kimi2026k3, zai2026glm53, qiu2026qwennext, xiaomi2026mimov26}. In most of their layers, the feed-forward network is split into many smaller networks, called \emph{experts}. A small network called the \emph{router} sends each \emph{token}, a word or a piece of a word, to a few of the experts. Only these experts process it. An MoE model can therefore have many parameters while each token uses only a small fraction of them. This makes it much cheaper to train and to run than a dense model of the same size, which uses all its parameters for every token.

Two choices in how these models are trained are the subject of this thesis. The first is how the experts are kept equally busy, which does not happen on its own. A router trained without any correction tends to send most tokens to a few experts. These experts are then trained on more tokens and become even more attractive to the router, while the others stay idle \cite{shazeer2017outrageously}. Many current models prevent this with a small correction for each expert, called the \emph{bias}. After every training step, the bias is lowered if the expert received too many tokens and raised if it received too few \cite{wang2024lossfree}.

For each token, the router computes a score for every expert, adds the bias of each expert to its score, and sends the token to the few experts with the highest sum. An expert that receives high scores for most tokens therefore needs a lower bias than the others, so that it receives no more than its \emph{fair share} of the tokens. The fair share is what each expert would receive if all of them were equally busy. The more strongly the router favors some experts over others, the further apart their biases must be. We measure how hard the bias works by the \emph{span}, the largest minus the smallest bias in a layer.

The second choice is the \emph{optimizer}, the rule that turns the gradient of the loss into an update of the parameters. Until recently, most language models were trained with AdamW \cite{loshchilov2019decoupled}, a variant of the Adam optimizer \cite{kingma2015adam} that scales the update of each parameter separately from the others. Most of the leading open models of 2026 instead train nearly all their weight matrices with Muon \cite{jordan2024muon, liu2025moonlight} (\cref{tab:related-router}), which treats each matrix as a whole and gives it an update of a fixed size. Since the optimizer trains the router and the rest of the network, it shapes the scores that the bias has to correct. This thesis therefore studies how the choice of optimizer changes the work of the bias.

The labs that build these models disagree on which optimizer should train the router. DeepSeek-V4 and Kimi K3 train it with Muon, like their other weight matrices \cite{deepseek2026v4, kimi2026k3}, whereas Qwen3.8-Flash-Next and MiMo-V2.6 use AdamW for it \cite{qiu2026qwennext, xiaomi2026mimov26}. According to the Qwen3.8-Flash-Next report, Muon made the router unstable early in training and gave no clear gain later \cite{qiu2026qwennext}. The DeepSeek-V4 report also links the routing to the stability of training. The \emph{loss spikes} of DeepSeek-V4, sudden jumps of the loss, came with very large activations in the MoE layers, which the routing itself appears to make larger \cite{deepseek2026v4}. None of these reports compare the options under controlled conditions. Moreover, the few studies that compare optimizers in MoE models either train without the bias or do not measure what it does \cite{wang2026depth, lau2026symmetry, srivastava2026lossfree}. Therefore, no published measurement shows how Muon changes the work of the bias or which part of the network causes this change.

\section{Objectives and Approach}
\label{sec:introduction-objectives}
The objective of this thesis is to measure how the optimizer changes the work of the bias and the behavior of the router in MoE language models. We pursue it through three research questions.
\begin{itemize}
  \item \textbf{RQ1.} How does training with Muon instead of AdamW change the span of the bias, the routing, the loss, and how specialized the experts become?
  \item \textbf{RQ2.} Current models use one of three \emph{partitions}, or assignments of the optimizers: all parameters on AdamW, all on Muon, or all on Muon except the router. How do these three partitions behave when the learning rate is raised far above its tuned value, as in the stress test of the Qwen3.8-Flash-Next report \cite{qiu2026qwennext}?
  \item \textbf{RQ3.} Which group of parameters, when trained with Muon, causes the higher span that RQ1 finds, and through what mechanism?
\end{itemize}

To answer these questions, we trained small MoE language models of two sizes on educational web text. The smaller uses 55M parameters for each token out of 303M in total, and the larger 108M out of 496M\@. In each run, we chose which optimizer trained each group of parameters, so that the router, the experts, and the other parts of the network could be moved to Muon one at a time. In total, we trained 145 runs. Except for the tuning of the learning rates, every configuration had three random seeds. Runs with the same seed start from the same weights and see the same data, so we compare configurations seed by seed. We also fixed the comparisons of each experiment before training its runs.

\section{Findings and Contributions}
\label{sec:introduction-findings}
Muon reached a lower loss than AdamW and made the experts more specialized, but it also made the bias work harder. The span was higher at every seed: 0.137 against 0.105 averaged over the layers, and 0.194 against 0.084 in the first MoE layer, where the difference was largest. We had expected the opposite. With Muon, the size of the update of an expert does not depend on how many tokens it receives, which should weaken the tendency of busy experts to become busier. Our runs confirmed this independence, yet the span still rose. How much it rose depended on the learning rate of Muon, which our tuning had set to twice that of AdamW\@. At half this rate, which trained slightly better, about three quarters of the extra span of the first MoE layer disappeared. The span also stayed higher with a bias four times faster, with the balance loss of DeepSeek-V3, and on the larger model.

The optimizer of the router changed the span of the first MoE layer little. What mattered was the optimizer of the \emph{backbone}, which we define as every weight matrix except the routers, the embedding table, and the output head. When the whole backbone was trained with Muon, the inputs of the first router shared a direction. The mean cosine similarity between the inputs of two tokens was about 0.4, against 0.07 under AdamW\@. Through this direction, the router favored the same experts for every token, and the bias needed hundreds of steps to cancel this preference. Meanwhile, the busiest expert of the layer received four to seven times its fair share. This overload did not raise our loss, but it could slow down training when the experts sit on many devices. In an exploratory test of whether the shared direction causes the overload, we repeated the start of some runs with one change. We subtracted the mean input of the batch from the input of the first router, which removes the shared direction. The overload and most of the extra span disappeared, and the loss did not change. Muon on the routed experts alone, however, raised the span even more, without a strong shared direction, and we do not know how.

A router trained with AdamW was updated more along the shared direction than one trained with Muon. With this router, the loads shifted more during the warmup, and the loss was slightly higher. This does not support the reason that the Qwen3.8-Flash-Next report gives for putting the router on AdamW\@. Far above the tuned rates, no run diverged, but the loss of Muon rose more than that of AdamW.

Summing up, the contributions of this thesis are the following:
\begin{itemize}
  \item A paired measurement, at two model sizes, of how much harder the bias works under Muon than under AdamW.
  \item A breakdown of this effect by group of parameters and by learning rate.
  \item An explanation of the higher span of the first MoE layer, tested by removing the shared direction.
  \item Evidence on which optimizer should train the router when the rest of the network is trained with Muon.
  \item A stress test of the three partitions of RQ2 at up to 16 times their tuned learning rates.
\end{itemize}

\section{Structure of the Document}
\label{sec:introduction-structure}
\Cref{sec:background} presents the background, and \cref{sec:related-work} reviews the related work. \Cref{sec:method} describes the setup that the experiments share. \Cref{sec:results} reports the experiments, and \cref{sec:discussion} interprets them and states their limits. \Cref{sec:conclusions} summarizes the answers and proposes future work. \Crefrange{app:configuration}{app:reproduction} give every setting of the runs, additional results, and how to reproduce them.
